%================================================================
% CHAPTER 6: CONCLUSION AND FUTURE WORK
%================================================================
\chapter{CONCLUSION AND FUTURE WORK}

\section{Introduction}

This chapter consists of the conclusion, future work and limitations. In this section we have explained in detail the conclusion of the whole project.

\section{Conclusion}

Academia Portal is designed for students and persons related to the educational field. This system is designed to help students overcome the issues related to the educational system. It is designed to facilitate students to get access to different university information that is available on the internet, all at one platform, and also to overcome the issues of manual search.

Finally, we conclude that the proposed system provides one platform on the internet. This system has the following qualities and characteristics:

\begin{itemize}
    \item It is loaded with a primitive menu.
    \item This system is capable of extracting data from different university websites automatically, related to general information like fee structure, programs offered, university criteria, faculty, ranking, and comparisons which help students related to their educational career.
    \item It is designed to be user-friendly where users can easily get their desired results by using specific functions.
    \item This system is developed using different updated algorithms like web scraping, web crawling and filtration, which will not only extract the data or contents but also filter the specified data before storing it in the database.
    \item Different functionalities of Academia Portal are data extraction, filtration, data crawling and data storage.
\end{itemize}

\section{Limitations}

This project is not accessible on all domains of the internet. The main target of this system is to access universities in the twin cities (Islamabad/Rawalpindi), which is our main limitation. This system is only specified for websites with extensions ``\texttt{.pk}'' and ``\texttt{.com}''. This system is limited because it accesses only some universities from the twin cities (Islamabad/Rawalpindi).

\section{Future Work}

In the existing system at this time it is for twin city (Islamabad/Rawalpindi) universities, but in the future it will be extended more. Then it will be extended to cover all universities present in Pakistan. Now this system is basically working for a countable number of universities, but in future after doing more work it will be further polished and made easy for all types of users.

Now, in this system some modules are included such as ranking, fee comparison, fee structure, eligibility criteria, faculty, and program offered, but some more modules will be added in the system to fulfill the needs of users. In the future this project will be extended to the country level, covering all provinces and ultimately all universities in Pakistan.
